\documentclass{article}
\usepackage[T1]{fontenc}
\usepackage{lmodern}
\usepackage{graphicx}
\usepackage{amsmath,amssymb}
\usepackage[a4paper,margin=2cm]{geometry}
\usepackage[absolute,overlay]{textpos}
\setlength{\TPHorizModule}{1mm}
\setlength{\TPVertModule}{1mm}
\usepackage[indonesian]{babel}
\usepackage{hyperref}
\setlength{\emergencystretch}{2em}
% unit: Halaman 14

\begin{document}

% === Halaman 14 (python/native) ===

Rinaldi M/II4020 Kriptografi

14

• Eve (seorang kriptanalis) yang menyadap pembicaraan Alice dan Bob tidak dapat 

menghitung K. 

• Sebab Eve hanya memiliki informasi p, g, A dan B yang semuanya tidak rahasia, 

tetapi ia tidak mengetahui nilai a atau b.  

• Untuk mengetahui a , Eve perlu melakukan perhitungan untuk menemukan a dari 

persamaan A = ga mod p.   (a adalah logaritma diskrit dari A dalam modulus p)

• Sekali a diketahui, maka selanjutnya Eve menggunakannya untuk menghitung 

kunci rahasia K = Ba mod p. 

• Kabar baiknya, logaritma diskrit bilangan bulat yang besar sangat sulit dihitung. 

Jadi, Eve tidak dapat menemukan kunci K.

\end{document}
